\section{The root eigenline and its deepest residue}
\label{hodge:section}

The coefficient-one part of the boundary naturally produces a mixed
variation.  We first identify the line of volumes inside that variation
and then identify the same line in a pure variation associated with a
deepest boundary stratum.  The comparison will retain the specified
volume, which is necessary for the extension calculation in
Section~\ref{orders:section}.

\subsection{A local root construction}
\label{hodge:rootconstruction}

Use the notation of Section~\ref{sec:prelim}.  Thus \(S\) is a smooth
compact K\"ahler modification with good open \(U\), the complement
\(S\setminus U\) is an SNC divisor, and
\[
 p:W\longrightarrow S,\qquad
 D_W=h^*\Delta-K_{W/Y},\qquad
 K_{W/S}+D_W\sim_{\Q}p^*Q_S.
\]
The equivalence is equipped with the actual bundle identification in a
fixed degree \(m\).  Over \(U\), the map and the boundary strata are
smooth, the boundary is relatively SNC with coefficients in \([0,1]\),
and
\[
 \mathcal G_m
 =p_*\cO_W\bigl(m(K_{W/S}+D_W)\bigr)|_U
 \simeq \cO_S(mQ_S)|_U.
\]
Only this good-open condition will be used in the present section.

Choose a sufficiently small coordinate neighborhood \(V\subset S\)
and a frame \(s\) of \(\cO_S(mQ_S)|_V\).  Its image on \(W_V\) is a
meromorphic section of \(\omega_{W/S}^{\otimes m}\), with divisor
\(-mD_W\).  Form its normalized \(m\)-th-root cover.  This construction
does not require a meromorphic trivialization of \(\omega_{W/S}\):
choose an effective integral divisor \(J\) that clears the poles, form
the finite root cover of the holomorphic section in
\((\omega_{W/S}(J))^{\otimes m}\), and normalize.  The resulting cover
with its meromorphic root is independent of the pole clearing.  In a
local canonical frame \(\alpha\), if \(s=a\alpha^{\otimes m}\), it is
given by adjoining \(z\) with \(z^m=a\), and the root form is
\(z\alpha\), pulled back as a differential form.

Keep every component of this cover.  The group \(\mu_m\) acts on it,
and the tautological root has a fixed character \(\chi\).  Take a
projective equivariant resolution and write
\[
 \rho:Z\longrightarrow W_V,\qquad
 \tau=\text{the pulled-back relative root form}.
\]
The space \(Z\) and its fibers are allowed to be disconnected.  The
finite-cover and resolution facts from Section~\ref{sec:prelim}
give K\"ahler total spaces after shrinking \(V\), with the resolution
performed over a slightly larger neighborhood.

Put \(V^\circ=V\cap U\).  On this open the construction can be made
smooth over the base, with a relative reduced SNC divisor \(H\)
consisting of the inverse images of the coefficient-one components.
Here is the local reason for this assertion.  In relative SNC
coordinates the tensor has the form
\begin{equation}
 s=
 a\,\prod_{\nu=1}^{r}x_\nu^{-m\delta_\nu}
 (dx_1\wedge\cdots\wedge dx_d)^{\otimes m},
 \qquad a\in\cO^*,\quad 0\leq\delta_\nu\leq1.
 \label{hodge:localtensor}
\end{equation}
After taking a local root of the unit and removing full \(m\)-th
powers, the normalized cover is a monomial cover in the coordinates
with \(0<\delta_\nu<1\).  The coefficient-one coordinates, the other
fiber coordinates, and the base coordinates are smooth product
parameters.  Functorial resolution, which commutes with smooth
morphisms, resolves the monomial part while preserving these product
parameters; one can equally use a toroidal resolution in the fractional
directions.  Thus the inverse images of the coefficient-one
components are smooth reduced divisors, their intersections are the
corresponding resolved covers of the original intersections, and all
these strata are smooth over \(V^\circ\).  This use of functorial
resolution is as in \cite[Theorems~1.1 and~1.3]{BM08}.

\begin{lemma}
\label{hodge:rootpoles}
Over \(V^\circ\), the form \(\tau\) is a relative top form with at most
logarithmic poles on \(H\) and no other poles.  At a strict covering
prime over a horizontal divisor of coefficient \(\delta\), if \(j\)
is the ramification index, its order is
\[
                    j(1-\delta)-1.
\]
\end{lemma}

\begin{proof}
At the generic point of the divisor, substitute \(x=z^j\) in
\(x^{-\delta}dx\).  The order is \(j-1-j\delta\).  It is integral;
for \(\delta<1\) it lies between \(0\) and \(j-1\), while for
\(\delta=1\) it equals \(-1\).

For exceptional divisors in a monomial resolution in the fractional
directions, the shifted order is the sum of the nonnegative coordinate
orders multiplied by \(1-\delta_\nu\).  At least one fractional
coordinate has positive order, so this sum is strictly positive.
The actual order is integral, hence nonnegative.  Equivalently, write
the root of \eqref{hodge:localtensor} in a logarithmic differential
basis.  The coefficient has a strictly positive gain in every
exceptional fractional direction, whereas the logarithmic Jacobian
has at most a simple pole.  The coefficient-one coordinates remain
transverse product parameters.  Consequently the only poles are the
asserted simple logarithmic poles along \(H\).
\end{proof}

\subsection{Extension facts and the mixed top step}
\label{hodge:extensionfacts}

We record precisely the extension input needed below.  Let
\(\Delta_t\) be a disk and \(\Delta_t^*=\Delta_t\setminus\{0\}\).
Suppose
\[
 g:(Z,H)\longrightarrow\Delta_t
\]
is proper, \(Z\) is smooth K\"ahler, \(H\) is a reduced SNC divisor,
every stratum of \((Z,H)\) dominates the disk, and all strata are
smooth over \(\Delta_t^*\).  Write \(d=\dim Z-1\).
Fujino--Fujisawa's Theorem~1.1 supplies the graded-polarizable real
variation on \(H_c^d(Z_t\setminus H_t)\) and its lower canonical
extension, with locally free double Hodge/weight grades
\cite[Theorem~1.1]{FF25}.  Its dual description gives
\begin{equation}
 g_*\omega_{Z/\Delta_t}(H)
 \simeq
 \left(\Gr_F^0\,
       \overline{H_c^d(Z_t\setminus H_t)}^{\,\mathrm{lower}}\right)^*.
 \label{hodge:FF}
\end{equation}
Poincar\'e duality identifies \(H_c^d{}^*\) with \(H^d(d)\).
Therefore \eqref{hodge:FF} is the extended \(F^d\) step of ordinary
cohomology in the \emph{upper} canonical extension.  All its double
grades remain locally free.  Upper residues lie in \((-1,0]\), and
lower residues in \([0,1)\); in the unipotent case the conventions
coincide \cite[Section~2.14]{FF25}.

For pure variations we will use the following consequences of the
nilpotent orbit theorem.  We include the functorial consequences in
the statement to specify exactly what is required of an extension.

\begin{lemma}[Pure extension rules]
\label{hodge:purerules}
Let a real-polarizable variation of pure Hodge structure be given on
the complement of an SNC divisor in a polydisk, with quasi-unipotent
local monodromies.
\begin{enumerate}[label=\textup{(\roman*)}]
\item
After coordinate power substitutions making all local monodromies
unipotent, its Hodge filtrations extend by subbundles of the canonical
flat bundle.  On a disk, the filtration also extends by subbundles
of either the upper or the lower canonical bundle before
unipotentization.
\item
After coordinate power substitutions making the monodromies
unipotent, the canonical extension and its Hodge filtration commute
with further powers and with pullback to any disk whose punctured
image lies in the given open.
\item
A real pure Hodge morphism has flat Hodge kernels and images, with flat
Hodge complements.  Its split maps and finite-group eigenspace
projectors are compatible with the canonical extensions in
\textup{(i)}.  After unipotentization, tensor constructions are
compatible as well.
\item
On a unipotent disk, a nonvanishing frame \(e\) of an extended highest
Hodge line satisfies, for some constants \(C,A>0\),
\begin{equation}
 C^{-1}(1+|\log|t||)^{-A}
 \leq \|e(t)\|
 \leq C(1+|\log|t||)^A.
 \label{hodge:logbounds}
\end{equation}
The constants may depend on the disk.  The same assertion applies to
Hodge determinant lines.
\end{enumerate}
\end{lemma}

\begin{proof}
In the unipotent case, extension by subbundles is the nilpotent orbit theorem
\cite[Theorem~4.12]{Schmid73}; a formulation with real polarization
and its canonical-frame interpretation is given in
\cite[Section~2, \((2.1)\)--\((2.2)\)]{CK89}.
For a quasi-unipotent disk, make the cyclic substitution \(t=u^N\)
that removes the semisimple monodromy.  Its finite deck group acts
on the extended filtered bundle.  Choose a character basis at the
origin adapted to the Hodge flag.  Lift these vectors to local
sections in the corresponding Hodge subbundles and average with
the character projectors.  The resulting equivariant sections
form an adapted frame near the origin.  Multiplying its vectors
by the integral powers of \(u\) prescribed by their characters
gives a descending frame with residue eigenvalues in either chosen
interval \((-1,0]\) or \([0,1)\).  The same subsets of this frame
span the respective extended Hodge steps.  This proves the disk
assertion for both canonical conventions.

For unipotent monodromies with commuting logarithms \(N_i\), the
canonical flat bundle is described in logarithmic frames.  Under a
disk map, each boundary coordinate is \(t^{a_i}v_i(t)\), with \(v_i\)
a unit and \(a_i\geq0\).  The pulled-back logarithmic connection has
nilpotent residue \(\sum_i a_iN_i\).  Thus the pulled-back flat bundle
is the canonical one.  The pulled-back Hodge subbundles are the
canonical filtration as well, by its characterization inside this
bundle.

For a pure Hodge map, its kernel and image are flat real Hodge
subbundles.  The restriction of a polarization to a real Hodge
subspace is nondegenerate.  Orthogonal complements therefore split
the map in the category of real variations.  These projectors are
flat, so in logarithmic frames their extensions have the same
constant ranks.  A polarization can be averaged over a finite group;
one then complexifies and applies the character projectors.
This uses no semisimplicity assertion for arbitrary local systems.

The abstract norm estimate is
\cite[Theorem~\(6.6'\)]{Schmid73}.  Applied to a flat frame and its
dual, it bounds the metric and its inverse by powers of
\(1+|\log|t||\).  Logarithmic frames differ from flat frames by finite
polynomials in the logarithm.  A nonvanishing extended line frame
has bounded coefficients in a logarithmic frame and one coefficient
bounded away from zero locally.  The metric and inverse-metric bounds
give \eqref{hodge:logbounds}.  Tensor and determinant constructions
preserve this conclusion.
\end{proof}

The variations occurring here have integral lattices up to isogeny.
Indeed they arise from the cohomology of compact SNC strata and the
rational weight filtration.  A K\"ahler class on the total space
restricts to flat real K\"ahler classes on these strata, so their pure
cohomologies, and the relevant kernels and images, are real-polarizable.
Their local monodromies are quasi-unipotent by the integral monodromy
theorem.  One may use its real-polarized formulation; alternatively,
the elementary transport in Lemma~\ref{periods:rationalization}
reduces it to the rationally polarized statement without changing
the integral monodromy.  Applying the theorem to weight grades also
gives quasi-unipotence for the mixed variation.

\begin{lemma}[The top step and its weight grade]
\label{hodge:doublegrades}
Let a mixed variation on a punctured disk be extended in its upper
canonical flat bundle, with the canonical weight filtration and
locally free double Hodge/weight grades, as in \eqref{hodge:FF}.
Suppose a finite-group character summand satisfies
\[
 \dim F^d=1,\qquad F^{d+1}=0,
\]
and the sole nonzero \(F^d\) weight grade has weight \(w\).
Then projection to \(\Gr^W_w\) identifies the extended top line with
the canonically extended top line of that pure grade.
\end{lemma}

\begin{proof}
The induced filtration of \(F^d\) by the weights has just one
nonzero successive quotient.  Local freeness of the double grades
implies that the same ranks and exact quotient sequences hold at the
origin.  Its terms below weight \(w\) vanish and its terms from
weight \(w\) onward equal the whole line.  Projection is therefore
an isomorphism there as well.  The upper canonical extension is
exact on the weight filtration, and the induced filtration on a
pure grade is its canonical Hodge filtration.  These statements also
follow by dualizing the lower extension in \eqref{hodge:FF}.
The character projector is an idempotent on all these locally free
objects, so taking its image preserves the argument.
\end{proof}

\subsection{The pure line detected by residues}
\label{hodge:pureresidue}

Over \(V^\circ\), put
\[
 \mathbb V=R^d(g|_{Z\setminus H})_*\C,
 \qquad g:Z\longrightarrow V,
\]
where the notation denotes ordinary cohomology on the smooth open
fibers.  We use \(W_\bullet\) for its weight filtration.  The
logarithmic mixed Hodge construction identifies
\[
 F^d\mathbb V_x
 \simeq H^0(Z_x,\omega_{Z_x}(H_x)).
 \tag{\(\ast\)}
\]
It is the ordinary-cohomology version of the compact-support SNC
construction above: compact supports are computed by the simple
complex of the compactification and its successive intersections;
duality gives logarithmic forms and residues.  These constructions
work for compact K\"ahler strata, using their pure Hodge structures
and Hodge degeneration, as in \cite[Section~4]{FF25}; compare
\cite[Section~3]{Deligne71}.

\Needspace{9\baselineskip}
\begin{proposition}[Root eigenline and deepest residue]
\label{hodge:residue}
Let \(k\) be the maximal number of horizontal coefficient-one
components meeting in a good fiber, constant after shrinking \(U\).
For the local root family above:
\begin{enumerate}[label=\textup{(\roman*)}]
\item
\[
                   F^d\mathbb V_\chi=\cO_{V^\circ}\tau.
\]
This line lies in the single pure weight grade
\(\Gr^W_{d+k}\mathbb V_\chi\).
\item
Choose a depth-\(k\) intersection component downstairs dominating
the good open.  After a good-open base change selecting one
connected component \(A_x^*\) of its fibers, take every root-cover
sheet above that component.  Iterated residue identifies the line
in \textup{(i)}, through its pure grade and with the Tate twist
removed, with the highest eigenline of the resolved root cover of
the induced klt volume on \(A_x^*\).  In degree \(m\) this is the
specified plurivolume identification from SNC adjunction.
\item
These highest-line identifications preserve generators in
canonical extensions on unipotent disks.  In a disk model satisfying
the hypotheses of \eqref{hodge:FF}, its extended mixed top line
projects isomorphically to the same pure line.
\end{enumerate}
If \(k=0\), the intersection is the whole original fiber, there is
no removed divisor, and the pure root-cover cohomology itself
provides the line.
\end{proposition}

\begin{proof}
By Lemma~\ref{hodge:rootpoles}, \(\tau_x\) belongs to the top
logarithmic-form space.  If \(\eta\) is another such eigenform with
character \(\chi\), the ratio \(\eta/\tau_x\) is invariant under
\(\mu_m\).  It descends to a meromorphic function on the original
connected compact smooth fiber, also when the normalized cover is
disconnected: the invariant meromorphic algebra of the full root
cover is the downstairs algebra.  At a strict covering divisor, a
downstairs pole of order \(n\geq1\) would lower the order in
Lemma~\ref{hodge:rootpoles} by \(jn\).  For \(\delta<1\) this gives a
pole where no pole is allowed; for \(\delta=1\) it gives an order
below \(-1\).  The ratio has no pole along any downstairs prime.
It extends across the remaining codimension-two set, and compactness
makes it constant.  This proves the first displayed equality.

There is no Hodge step above \(F^d\).  Exactness of the induced
weight filtration on this one-dimensional space shows that exactly
one weight grade carries it.  To determine that grade, let
\(H_x^{[r]}\) be the disjoint union of the \(r\)-fold intersections,
with \(H_x^{[0]}=Z_x\).  The residue spectral sequence has terms
\begin{equation}
 E_1^{-r,q}
   =H^{q-2r}(H_x^{[r]},\C)(-r)
 \ \Longrightarrow\ H^{q-r}(Z_x\setminus H_x,\C).
 \label{hodge:weightspectral}
\end{equation}
The differentials on the first page are pure Hodge maps.  Taking
their kernels modulo images leaves pure objects, and subsequent
differentials vanish by their different weights.  Since
\(H_x^{[k+1]}\) is empty, there is no incoming first differential at
\((-k,d+k)\).  Consequently \eqref{hodge:weightspectral} gives a
natural pure map
\begin{equation}
 \Gr^W_{d+k}H^d(Z_x\setminus H_x,\C)
 \lhook\joinrel\longrightarrow
 H^{d-k}(H_x^{[k]},\C)(-k).
 \label{hodge:residuemap}
\end{equation}
On the highest Hodge step it is iterated residue.  In local SNC
coordinates this extracts the coefficient of
\(dx_1/x_1\wedge\cdots\wedge dx_k/x_k\); integration over the
corresponding normal circles gives the same cohomological map,
up to fixed conventional constants.  Thus it is horizontal in
families.

At a point of a selected deepest intersection away from fractional
boundary, the root cover is unramified and \(\tau_x\) has exact
logarithmic poles in the \(k\) transverse coordinates.  Its residue
is nonzero.  It extends as a holomorphic top form on the compact
resolved intersection upstairs and hence has a nonzero cohomology
class.  The unique weight carrying \(F^d\mathbb V_\chi\) is therefore
\(d+k\).

Consider now the selected connected component \(A_x^*\) downstairs.
By maximality of \(k\), no further coefficient-one component meets
it.  The remaining boundary is effective fractional SNC, so the
induced pair is klt.  In the product coordinates of the root
construction, taking residue simply deletes the \(k\) logarithmic
parameters.  The union of all sheets above \(A_x^*\) is thus the
resolved root cover of its induced log plurivolume.  Its top
eigenline has dimension one by the same ratio argument.  The
projection of \eqref{hodge:residuemap} to these sheets is nonzero on
the source line, and therefore identifies the two top lines.
With an ordering of the boundary components, tensor power of
ordinary residue gives
\[
       (\operatorname{Res}\tau_x)^{\otimes m}
       =\rho^*(\operatorname{Res}_m s_x).
\]
Here \(\operatorname{Res}_m\) denotes the bundle map supplied by
SNC adjunction; fixed sign or normalization choices do not affect
the argument.

The maps on pure grades and compact intersection cohomology are
real Hodge maps before taking the character summand.  Their
kernels and images split by Lemma~\ref{hodge:purerules}, including
after projection to all the selected sheets.  They hence preserve
the extended highest-line generators.  On a disk to which
\eqref{hodge:FF} applies, Lemma~\ref{hodge:doublegrades} identifies
the mixed top line with that on its unique pure grade.  This proves
\textup{(iii)}.

For \(k=0\) there is no logarithmic boundary, so ordinary compact
root-cover cohomology is pure of weight \(d\), and the same
eigenform argument applies directly.
\end{proof}

\subsection{Changing the resolved root model}
\label{hodge:modelchanges}

The order calculation will use resolved models adapted to a chosen
disk.  Such a model can modify a generic root fiber, so one needs
the following invariance statement rather than a claim that every
cohomology group is birationally invariant.

\begin{lemma}
\label{hodge:birational}
Over a smooth good open, let
\(\varphi:(Z',H')\to(Z,H)\) be a proper equivariant modification
between smooth families of SNC pairs, with \(H'\) the reduced
inverse image of \(H\).  Pullback is an isomorphism on top
logarithmic-form spaces.  In the character summand of
Proposition~\ref{hodge:residue}, it identifies the top line and
the top line of its unique pure weight grade.

Suppose further that, over a disk, the new compactification
satisfies the hypotheses of \eqref{hodge:FF}.  If the original
variation has been made unipotent, testing extension of the
pulled-back volume in \(g'_*\omega_{Z'/\Delta}(H')\) tests the
canonical extension of the original top line.  It is not necessary
to assume unipotence of every summand newly appearing in the
cohomology of \(Z'\setminus H'\).
\end{lemma}

\begin{proof}
Pullback of an SNC logarithmic top form has at most logarithmic
poles along the reduced inverse image of the boundary.  To descend
a top logarithmic form, use the isomorphism away from the
codimension-two exceptional image on each smooth original fiber.
The resulting section of \(\omega_Z(H)\) extends across that image
because this is a line bundle.  Thus pullback and descent are
inverse on top forms.

Pullback of logarithmic complexes gives the cohomological morphism
of mixed variations.  It respects residues and both filtrations;
one can see the logarithmic pullback in SNC coordinates, and the
compact-stratum construction gives the same rational map.
Morphisms of mixed Hodge structures are strict for the weight
filtration on each Hodge step.  Since the top line maps
isomorphically, its unique carrying weight is unchanged, and the
induced pure map is an isomorphism on these highest lines.

Use the upper canonical extension for the new model.  Its extended
mixed top eigenspace projects isomorphically to its pure grade by
Lemma~\ref{hodge:doublegrades}.  The comparing pure map splits onto
its image and is compatible with upper extension by
Lemma~\ref{hodge:purerules}.  That image is the unipotent summand
pulled from the original variation, so upper extension on it is
the unipotent canonical extension.  All other weight grades have
zero top step of this character.  Formula~\eqref{hodge:FF} therefore
tests precisely the original extended line, as claimed.
\end{proof}
